\chapter*{Introduction}
\phantomsection
\addcontentsline{toc}{chapter}{Introduction}

The purpose of this Book is the elementary study of the infinitesimal properties of one real variable; the extension of these properties to functions of several real variables, or, all the more, to functions defined on more general spaces, will be treated only in later Books.

The results which we shall demonstrate will be useful above all in relation to (finite) real-valued functions of a real variable; but most of them extend without further argument to functions of a real variable taking values in a topological vector space over R (see below); as these functions occur frequently in Analysis we shall state for them all results which are not specific to real-valued functions.

The notion of a topological vector space, of which we have just spoken, is defined and studied in detail in Book V of this Series; but we do not need any of the results of Book V in this Book; some definitions, however, are needed, and we shall reproduce them below for the convenience of the reader.

We shall not repeat the definition of a vector space over a (commutative) field K (Alg., II, p.~193). \(^{1}\) A topological vector space E over a topological field K is a vector space over K endowed with a topology such that the functions \(x + y\) and xt are continuous on \(E \times E\) and \(E \times K\) respectively; in particular, such a topology is compatible with the structure of the additive group of E. All topological vector spaces considered in this Book are implicitly assumed to be Hausdorff. When the topological group E is complete one says that the topological vector space E is complete. Every normed vector space over a valued field K (Gen.~Top., IX, p.~169) \(^{2}\) is a topological vector space over K.

Let E be a vector space (with or without a topology) over the real field R; if x, y are arbitrary points in E the set of points \(\mathbf{x}t + \mathbf{y}(1 - t)\) where t runs through the closed segment {[}0, 1{]} of R is called the closed segment with endpoints x, y. One says that a subset A of E is convex if for any x, y in A the closed segment with endpoints x and y is contained in A. For example, an affine linear variety is convex; so is any closed segment; in \(R^{n}\) any parallelotope (Gen.~Top., VI, p.~34) is convex. Every intersection of convex sets is convex.

We say that a topological vector space E over the field R is locally convex if the origin (and thus any point of E) has a fundamental system of convex neighbourhoods. Every normed space is locally convex; indeed, the balls \(\|x\| \leqslant r (r > 0)\) form a fundamental system of neighbourhoods of 0 in E, and each of these is convex, for the relations \(\|x\| \leqslant r\) , \(\|y\| \leqslant r\) imply that

\[
\| \mathbf {x} t + \mathbf {y} (1 - t) \| \leqslant \| \mathbf {x} \| t + \| \mathbf {y} \| (1 - t) \leqslant r
\]

for \(0 \leqslant t \leqslant 1\) .

Finally, a topological algebra A over a (commutative) topological field K is an algebra over K endowed with a topology for which the functions \(x + y\) , xy and xt are continuous on \(A \times A\) , \(A \times A\) and \(A \times K\) respectively; when one endows A only with its topology and vector space structure over K then A is a topological vector space. Every normed algebra over a valued field K (Gen.~Top., IX, p.~175) is a topological algebra over K.

